\section{Evaluation Protocol}
\enlargethispage{2\baselineskip}

All experiments in this thesis are conducted on a single backbone model, Qwen2.5-7B. The evaluated baselines are FT, ROME, MEMIT, GRACE, WISE, and AlphaEdit. The study intentionally fixes the model family so that differences in the reported behavior can be attributed to the editing methods themselves rather than to backbone, tokenizer, or prompt-format differences.

Edits are applied sequentially and the edited model is evaluated at selected cumulative milestones. In principle, the protocol supports a much broader range of checkpoints, including substantially larger milestones as the edit stream grows. In the current report, the main analysis emphasizes milestones such as 1, 5, 10, 50, and 100 because they offer a practical balance between runtime and the need to observe when degradation begins to appear. At a given milestone $k$, the reported value is computed from the model state after the first $k$ edits in the stream. This setup reflects a lifelong protocol rather than a one-shot or purely batched setting.

Both cloze-style and QnA-style prompts were prepared in the evaluation data. In the present chapter, the main reported tables use the cloze track, while the QnA track is used only as an auxiliary comparison for prompt-format sensitivity. In either track, model outputs are compared against the target answer under the metric definitions given in the previous section. The scoring procedure is intended to accept coherent correct outputs while rejecting corrupted or repetitive generations. This distinction is important in the present study because some methods, especially under unstable sequential updates, can produce generations that contain the target token but are still unusable as answers.

An important interpretation detail is that the milestone results are cumulative evaluations on growing prefixes of the edit stream. Later milestones are therefore not computed on exactly the same prompt set as earlier milestones. This means that changes in a score can reflect both the model's cumulative editing behavior and changes in the composition of the evaluated prefix. The resulting curves should be read as comparative evidence under a shared sequential protocol rather than as measurements on one fixed held-out set.
